\section{Why Public Security Data Cannot (Yet) Test Investigation Policies}
\label{sec:realdata}

A controlled environment invites the obvious objection: would any of this hold on real data? We tried to answer it with the three public datasets that come closest. Each measurement below used only training or development material. No test split was read, and the same pre-registration discipline applied. None of the three can test an investigation policy, and the reasons are informative.

\paragraph{What a usable dataset needs.} Testing whether a controller chooses probes well requires three things. (i) The alert alone must not reveal the answer: several explanations must remain plausible. (ii) Probes must differ in cost and in how much they reveal, so that the order of probing matters. (iii) Ground truth must be objective, reflecting what actually happened rather than a labeling convention.

\paragraph{ExCyTIn-Bench~\citep{excytin}: the question names the answer.} ExCyTIn poses questions over a simulated Azure tenant's logs. Its answers can be recovered mechanically from the published investigation graphs (98.5\% of 1{,}017 questions), and candidate hypothesis sets can be derived from them. Two measurements on the 418 training questions ruled it out. First, 54.8\% of questions name the target alert in the question text, and most others paraphrase it, so the task is to locate the named alert and read an entity: navigation and retrieval. Second, we tested the standard indicator pivot probe, ``find rows in table $T$ that mention a known indicator''. The true answer appears no more often than distractors do (e.g.\ 0.571 vs.\ 0.570 in the alert table), so the probes have no discriminative power. Only 19.4\% of questions have any table that singles out the answer.

\paragraph{GUIDE~\citep{guide}: the label is organizational policy.} GUIDE contains 707{,}108 real, anonymized incidents with triage grades (true positive, benign positive, false positive) from 5{,}340 organizations. We defined ten incident-level probes over its alert metadata. Organization identity alone carries 1.03 of the grade's 1.50 bits of entropy. Several large organizations assign one grade to every incident. Within an organization, looking up how that organization previously graded the same detector reaches 0.911 accuracy, and once the detector is known the other facets carry almost no information. Across organizations, the probes do not beat the majority class. On the 13.3\% of incidents that are cold starts for their organization, evidence integration looked at first like a recall gain. It turned out to be a shift of operating point. Measured by AUROC, the \hesp{}-style integration (0.755) equals history lookup (0.759), a difference of $-0.004$ $[-0.028,+0.027]$. We paused the planned GUIDE study, and its test split remains unread.

\paragraph{OTRF Security-Datasets~\citep{otrf}: the tools give themselves away.} OTRF contains lab recordings of Windows host telemetry captured while known ATT\&CK techniques were executed. All 113 recordings are attack recordings; there are no benign-activity recordings. We audited 27 recordings from the remote-execution family (243{,}000 events). We treated every service installation and scheduled-task creation or update as an alert and labeled it from the recording's own metadata. This gave 46 distinct alerts: 8 attacks and 38 background operating-system events. A one-glance rule on the alert text, which asks whether the task or service belongs to a Windows component, is correct on 44 of 46. Six of the eight attacks are named by their tooling, through random service names, lab-specific task names, known implant binaries, or encoded PowerShell. Only one scenario needs investigation, and a single probe resolves it.

\paragraph{The common gap.} The three datasets fail in three different ways, but the missing ingredient is the same. The samples that would test an investigation policy are attacks that look like routine administration and routine administration that looks like an attack. Public data lacks exactly these. Benign counterparts are operating-system housekeeping rather than administrators using the same tools legitimately, and attack emulations leave artifacts that no careful adversary would. We see building such paired, objectively labeled data as a prerequisite for evaluating triage agents on real telemetry.

\subsection{Case study: a scheduled task masquerading as a Windows task}
\label{sec:case}
The one OTRF scenario that does require investigation shows, on real logs, the kind of case \hesp{} is built for. The recording's metadata references Microsoft's analysis of the Solorigate intrusion~\citep{solorigate} and labels the activity as remote scheduled-task creation (T1053.005). All values below are extracted by a script from the raw recording.

\begin{enumerate}
  \item \textbf{Alert.} Security event 4698 on \texttt{WORKSTATION6}: a new task \path{\Microsoft\Windows\SoftwareProtectionPlatform\EventCacheManager}. The folder and naming style match genuine Windows tasks, and a real task from the same folder appears routinely in these recordings. The text-only rule calls it benign.
  \item \textbf{Probe: creator.} The task was created by the domain user \texttt{THESHIRE\textbackslash pgustavo}. The other 17 task events under \path{\Microsoft\Windows\} in the 27 recordings were all made by machine accounts. This contradicts the hypothesis that Windows itself created the task.
  \item \textbf{Probe: logon session.} The creating session is a network logon (type 3, Kerberos) from another host. This contradicts local persistence and a user-context updater, the only kind of user-created benign task in the data, and supports remote creation.
  \item \textbf{Probe: action and source.} The task runs \texttt{cmd.exe} as SYSTEM and names the user as its author. The source host's process log shows the \texttt{schtasks /create \dots{} /s} command that created it.
\end{enumerate}

Here the intuition suggested by the alert text and the conclusion reached by investigation are opposite, and two cheap probes separate them. The case is qualitative. Because one probe suffices, it cannot show that ordering matters, and it is a single scenario. We include it as motivation, not as evidence.
